% vim: textwidth=78 ai spell spelllang=en
\documentclass[11pt,a4paper]{article}

\usepackage[utf8x]{inputenc}
\usepackage[IL2]{fontenc}
\usepackage[final]{hyperref}

\setlength{\textwidth}{140mm}
\setlength{\textheight}{220mm}
\setlength{\oddsidemargin}{12mm}
\setlength{\topmargin}{8mm}

\font\logo fi-logo at 4cm

\setcounter{tocdepth}{2}

\addtolength{\oddsidemargin}{0.2in}
\addtolength{\evensidemargin}{0in}
\addtolength{\textwidth}{0in}
\addtolength{\topmargin}{0in}
\addtolength{\textheight}{0in}

\newcommand{\stan}{\textsc{Stanse}}

\begin{document}
\begin{titlepage}
\vspace{2cm}
\begin{center}
      {\sc Masaryk University\\Faculty of Informatics}

     \vspace{1cm}
     \logo SL

     \vspace{2cm}
     {\bf\huge Symbolic execution on OS kernels}\\[0.8in]

     {\sc\large Dissertation Thesis Topic}\\[0.4in]

     {\Large\bf Jiří Slabý}
     \vfill
     {\large Brno, spring 2010}
     \vfill
\end{center}
     {\bf Supervisor:} prof. RNDr. Antonín Kučera, Ph.D.\\
     {\bf Supervisor-specialist:} Mgr. Jan Obdržálek, PhD.\\[0.5cm]
     {\bf Supervisor's Signature: ......................}
\end{titlepage}

\pagestyle{empty}

\tableofcontents

\pagestyle{plain}

\newpage
\setcounter{page}{1}

\section{Introduction}
Static analysis is currently considered to be a powerful approach used for
finding errors in a source code. The first real application was performed by
Engler et al. in \textsc{Metal}\cite{metal} in 2000.

Static analysis might have a form of finding some patterns in the code in some
specified manner (similarly to \textsc{Metal}) or test cases generation which
try to cover the most of the code (covering both true and false branches of
each branching).

Many tools were developed based on either approach however they are commercial
with closed source code, limited in what they can do or both. We are currently
working under Institute of Theoretical Informatics at the Faculty of
Informatics on a free (open-source) tool called \stan. It aims at finding
errors with a use of static analysis.

\subsection{Symbolic Execution}
One method of static analysis, \emph{Symbolic execution} is currently widely
used for checking code properties. It stems from \cite{select} and many
studies tried to apply the principle on a variety of userspace programs.
\textsc{Klee}\cite{klee} primarily aims at tests generation that would cover
as many program instructions as possible. \textsc{Pex}\cite{pex} uses
symbolic execution in cooperation with standard execution on processor
to address program coverage as well, but from another point of view.
Section \ref{current-se} is deserved for detailed description of the work done
in this area.

To the best of my knowledge nobody has ever tried to apply symbolic execution
on an operation system kernel (further only \emph{kernel}). This task
incorporates not only the execution
itself, but some kind of code preparation, compilation, linking and dynamic
linking. These terms will be discussed further in the text.

\subsection{Justification}
Symbolic execution has shown its power in finding errors in userspace
programs. The idea of applying the approach to the kernel is straightforward
extension, however it's not as easy task as it might seem. The kernel needs to
be processed, the result then executed and checked for some properties.
All of
\begin{itemize}
\item processing the kernel code,
\item finding the right properties to check and
\item finding them in the processed result
\end{itemize}
was not done by anybody yet. The dissertation should lead to complete
implementation of the mentioned and finding errors with minimum false
positives.

Dissertation Thesis Intent; Methods of Solving the Problem; Comparison of Key
Results

\newpage

\section{Report on Current Results}
\subsection{Static analysis}
Static analysis is an old approach, first summarized and formalized by Cousot
\& Cousot in \cite{cousot}) in 1977.

which
came into people's mind even after 2000. It was thanks to work of Engler et
al. on tool called \textsc{Metal}\cite{metal}.

\textsc{Coverity Prevent}\cite{coverity}, \textsc{FindBugs}\cite{findbugs},
\textsc{Klocwork}\cite{klocwork}

\subsection{Symbolic Execution}
\label{current-se}
\textsc{Klee}\cite{klee}, \textsc{Pex}\cite{pex}

\newpage

\section{Dissertation Thesis Intent}
As mentioned in the introduction, the intention is to apply symbolic execution
on a kernel. Dissertation Thesis should contain all the steps needed to
perform the execution to find errors in the code. From processing the code
(compilation) towards recognizing specified error patterns and report them.
Now I'll go into details of each phase and describe the idea behind.

\subsection{Code processing}
Make build system support: \textsc{Sparse}\cite{sparse}

Compilation gnome: \textsc{LLVM}\cite{llvm}

Compilation problems: inlined assembly, 

Linking: \texttt{llvm-link}

Linking against unknown functions vs. userspace: \textsc{Klee}\cite{klee}'s
solution; functions written in assembler

\subsection{Specifying Error Patterns}
Instrumented code: asserts, \texttt{\_\_chk\_user\_ptr}

\subsection{Finding Error Patterns}
Marek's dissertation, fallback to \textsc{Klee} (we have source code).

Expected Results; Schedule for Further Development; Expected Outcomes

\newpage

\bibliographystyle{plain}

\bibliography{teze}
\addcontentsline{toc}{section}{References}

\newpage

\section{Summary of Study Results}
During the first year of doctoral study I studied papers with results of
past researches in the area of static analysis. Also I read through some
literature about static analysis and compilers.

Under \emph{ITI}, we developed \stan, a static analysis tool written in
Java.\\[5mm]\noindent
The following publications were accepted:
\begin{itemize}
\item J. Obdržálek, J. Slabý and M. Trtík. \emph{STANSE: Framework for
  Finding Bugs in C Programs}
\end{itemize}

\end{document}
